\documentclass[noauthor,noinstructornotes]{ximera}

\title{Exponent of Zero}
\author{Kenneth Berglund}

\begin{document}
\begin{abstract}
    Investigating what happens with a zero in the exponent.
\end{abstract}
\maketitle

Using the rules of exponents, answer the following questions.
\textbf{Please do not look ahead until you have finished the current part!}

\begin{enumerate}
    \item Solve the equation \(2^1 \cdot 2^{-1} = 2^x\).
    \item Simplify \(2^1\) and \(2^{-1}\). Then multiply your answers together.
    What is the result?
    \item To summarize, we've shown that \(2^1 \cdot 2^{-1} = 2^0\), but
    also that \(2^1 \cdot 2^{-1} = 1\). What can you say about the value
    of \(2^0\)?
    \item Was there anything special about the base of 2 in this
    problem? What is \(17^0\)?
    \item (If you have time and energy) The expression \(0^0\) has 
    different interpretations depending on who you ask. The argument
    we used above to show that \(2^0 = 1\) does not work for \(0^0\).
    Why not?
\end{enumerate}

\end{document}
